% 例 4 链条：∠1 与 ∠2 互余 (90°)，∠2 与 ∠3 互补 (180°)
\documentclass[tikz,border=4pt]{standalone}
\usepackage{ctex}
\usepackage{amsmath}
\usetikzlibrary{calc, arrows.meta}
\begin{document}
\begin{tikzpicture}[scale=1.0, thick]
  % 左：互余  ∠1=30, ∠2=60
  \begin{scope}[shift={(0,0)}]
    \coordinate (O) at (0,0);
    \draw[-{Stealth[length=3pt]}] (O) -- (2.5,0);
    \draw[-{Stealth[length=3pt]}] (O) -- (0,2.5);
    \draw[-{Stealth[length=3pt]}, red] (O) -- ({2.5*cos(30)}, {2.5*sin(30)});
    \draw (0.25,0) -- (0.25,0.25) -- (0,0.25);
    \draw (0.9,0) arc (0:30:0.9);
    \node[font=\footnotesize] at ({0.65*cos(15)}, {0.65*sin(15)}) {$\angle 1$};
    \draw ({0.9*cos(30)}, {0.9*sin(30)}) arc (30:90:0.9);
    \node[font=\footnotesize] at ({0.6*cos(60)}, {0.6*sin(60)}) {$\angle 2$};
    \filldraw (O) circle (1.2pt);
    \node[below, font=\small] at (1.2,-0.4) {$\angle 1+\angle 2=90°$};
  \end{scope}
  % 右：互补  ∠2=60, ∠3=120
  \begin{scope}[shift={(6,0)}]
    \coordinate (O) at (0,0);
    \draw[-{Stealth[length=3pt]}] (O) -- (2.5,0);
    \draw[-{Stealth[length=3pt]}] (O) -- (-2.5,0);
    \draw[-{Stealth[length=3pt]}, red] (O) -- ({2.5*cos(60)}, {2.5*sin(60)});
    \draw (0.9,0) arc (0:60:0.9);
    \node[font=\footnotesize] at ({0.6*cos(30)}, {0.6*sin(30)}) {$\angle 2$};
    \draw ({0.9*cos(60)}, {0.9*sin(60)}) arc (60:180:0.9);
    \node[font=\footnotesize] at ({0.6*cos(120)}, {0.6*sin(120)}) {$\angle 3$};
    \filldraw (O) circle (1.2pt);
    \node[below, font=\small] at (0,-0.4) {$\angle 2+\angle 3=180°$};
  \end{scope}
  % 中间 → 箭头
  \draw[->, thick] (3,1) -- (5,1);
  \node[above, font=\small] at (4,1) {链条};
\end{tikzpicture}
\end{document}
